%%%%%%%%%%%%%%%%%%%%%%%%%
%% Header for standard beamer presentation
%%
%%  PresentationHeader.tex
%%
%%%%%%%%%%%%%%%%%%%%%%%%%

\documentclass[english,10pt]{beamer}

%%%%%%%%%%%%%%%%%%%%
%% Include general header where common packages are defined
%%%%%%%%%%%%%%%%%%%%

\input{GeneralHeader.tex}

%%%%%%%%%%%%%%%%%%%%
%% Idem general commands
%%%%%%%%%%%%%%%%%%%%

%\input{/Users/Juste/Documents/ComplexSystems/CityNetwork/Docs/Headers/GeneralCommands.tex}


\usetheme{Boadilla}

%\setbeamertemplate{footline}[text line]{}
%\setbeamercolor{structure}{fg=purple!50!blue, bg=purple!50!blue}

% cybergeo palette (from Clem server)
% #1C6F91, "#df691a", "#77c5ba", "orange", "#2db92d", "#e1ff2f", "#ff2313", "#bbab61
% redefine palette
\definecolor{cybblue}{HTML}{1C6F91}
\definecolor{cyborange}{HTML}{DF691A}
\definecolor{cybbluegreen}{HTML}{77C5BA}
\definecolor{cybgreen}{HTML}{2DB92D}
\definecolor{cybgreenyellow}{HTML}{E1FF2F}
\definecolor{cybred}{HTML}{FF2313}
\definecolor{cybglaucous}{HTML}{BBAB61}


\setbeamercolor{structure}{fg=cybblue}
\setbeamercolor{footline}{fg=orange, bg=orange}

\setbeamercovered{transparent}




\addtobeamertemplate{frametitle}{\hspace{-0.4cm}\vspace{-0.1cm}\includegraphics[height=1.2cm,width=1.2\textwidth]{template/bandeau3}\\
}{%
%\begin{textblock*}{150mm}(-1cm,-1.5cm)
%\end{textblock*}
}



% shortened command for a justified frame
\newcommand{\jframe}[2]{\frame{\frametitle{#1}\justify{#2}}}


\newcommand{\noun}[1]{\textsc{#1}}
\newcommand{\jitem}[1]{\item \begin{justify} #1 \end{justify} \vfill{}}
\newcommand{\sframe}[2]{\frame{\frametitle{#1} #2}}

\DeclareMathOperator{\Cov}{Cov}
\DeclareMathOperator{\Var}{Var}
\DeclareMathOperator{\E}{\mathbb{E}}
\DeclareMathOperator{\Proba}{\mathbb{P}}

\newcommand{\Covb}[2]{\ensuremath{\Cov\!\left[#1,#2\right]}}
\newcommand{\Eb}[1]{\ensuremath{\E\!\left[#1\right]}}
\newcommand{\Pb}[1]{\ensuremath{\Proba\!\left[#1\right]}}
\newcommand{\Varb}[1]{\ensuremath{\Var\!\left[#1\right]}}

% norm
\newcommand{\norm}[1]{\| #1 \|}

\newcommand{\indep}{\rotatebox[origin=c]{90}{$\models$}}


\usepackage{textpos}



%%%%%%%%%%%%%%%%%%%%%
%% Begin doc
%%%%%%%%%%%%%%%%%%%%%

\begin{document}



\title{Cybergeo : Bibliom{\'e}trie indirecte par analyse de r{\'e}seau}


%\author{J.~Raimbault$^{1,2}$}

%\institute{$^{1}$G{\'e}ographie-cit{\'e}s (UMR 8504 CNRS)\\
%$^{2}$LVMT (UMR-T 9403 IFSTTAR)}


\date{R{\'e}union 27/01/2016}


\usebackgroundtemplate{%
  \AddToShipoutPicture*{
    \put(0,0){
        \parbox[b][\paperheight]{\paperwidth}{%
            \vfill
            \centering
            \includegraphics[width=\paperwidth]{figures/nw}%
            \vfill
        }
    }
    \put(0,0){%
        \transparent{0.9}\textcolor{white}{\rule{\paperwidth}{\paperheight}}
    }
}
  %\includegraphics[width=1.5\paperwidth,height=1.5\paperheight]{figures/nw}
} 

%%%%%%%%%%%%%%%%%%%%%%%%%%%%%%%%
\begin{frame}
\titlepage
\end{frame}

%\begin{frame}
%\tableofcontents
%\end{frame}
%%%%%%%%%%%%%%%%%%%%%%%%%%%%%%%%


\usebackgroundtemplate{}


\section{Réseaux Sémantique}


\jframe{Collection des donn{\'e}es}{

\textbf{Retour sur le r{\'e}seau de citations.} apr{\`e}s raffinement, $\simeq$ 947 r{\'e}f{\'e}rences de cybergeo exploitables, sur $\simeq 1200$ ; certaines inexistantes, d'autres mal enregistr{\'e}es sur scholar $\rightarrow$ possibilit{\'e} de compl{\'e}ter {\`a} la main (long...).

\bigskip

\textbf{R{\'e}seau s{\'e}mantique.} Collection des r{\'e}sum{\'e}s/ann{\'e}es/auteurs/mots-cl{\'e}s pour les 400000 r{\'e}f{\'e}rences via l'API Mendeley
$\rightarrow$ $\sim 215000$ r{\'e}f{\'e}rences avec donn{\'e}es compl{\`e}tes.

\textbf{Statistiques}

\textit{Langues :} anglais 206607, francais 4109, espagnol 2029, allemand 892, portugais 891, n{\'e}erlandais 124, autres 182

\textit{Repartitions par ann{\'e}es : }
\centering

\includegraphics[width=0.3\textwidth,height=0.3\textheight]{figures/years}

}



\jframe{Extraction des mots-cl{\'e}s brute}{

\textit{Text-mining en python avec \texttt{nltk}}

\begin{itemize}
\item Detection de la langue par \textit{stop-words} % TODO language stats
\item Parsing/tokenizing/pos-tagging/stemming effectué différemment selon la langue :
\begin{itemize}
\item Anglais : pos-tagger intégré à \texttt{nltk}, combin{\'e} à un \emph{PorterStemmer}
\item Français ou autre : utilisation de \texttt{TreeTagger}
\end{itemize}
\item Selection des \emph{n-grams} potentiels (avec $1 \leq n \leq 4$) : anglais $\bigcap \{NN \cup VBG \cup JJ \}$ ;  français  $\bigcap \{NOM \cup ADJ\}$
\item Insertion en base pour extraction quasi-instantanée plus tard (10j $\rightarrow$ 5min)
\end{itemize}

}


\jframe{Estimation de la pertinence par bootstrap}{
Estimation exacte de la pertinence via la repartition statistique des co-occurrences (score de $\chi^2$) : \textit{termhood} d{\'e}finie comme
\[
t_i = \sum_{j\neq i}\frac{\left( M_{ij} - \sum_{k}M_{ik} \sum_{k} M_{jk}\right)^2}{\sum_{k}M_{ik} \sum_{k} M_{jk}}
\]


en $\Theta (\sum_i N_i^2)$ ($N_i$ taille des r{\'e}sum{\'e}s) : impossible sur un corpus o{\`u} $\sum_i N_i^2 \simeq N <N_i>^2 \simeq 4\cdot 10^6$

\bigskip


$\rightarrow$ Estimation par \textit{boootstrap} sur des {\'e}chantillons du corpus : moyenne de la termhood sur $B$ {\'e}chantillons de taille $C$, avec nombre de mots cl{\'e}s $K_L$


}


\jframe{Bootstrap : convergence de l'estimateur}{
\includegraphics[width=0.55\textwidth,height=0.6\textheight]{figures/kw50}
\includegraphics[width=0.55\textwidth,height=0.6\textheight]{figures/kw200}
}





\section{Suite de l'étude}


\jframe{Suite}{
\begin{itemize}
\item Bootstrap parallelis{\'e} sur l'ensemble du corpus
\item Construction du r{\'e}seau : co-occurrences dans les r{\'e}f{\'e}rences\\\textbf{Q : } utilisation des mots-cl{\'e}s des m{\'e}tadonn{\'e}es ? si oui filtrage sur fr{\'e}quence ? test avec/sans
\item D{\'e}tection de disciplines par communaut{\'e}s dans le r{\'e}seau s{\'e}mantique
\item Croisement des deux couches pour extraire positionnement et importance disciplinaire de Cyberg{\'e}o
\end{itemize}

}

\jframe{R{\'e}sultats attendus}{
\begin{itemize}
\item Couche des citations : analyse plus fines, cliques et communaut{\'e}s
\item Interdisciplinarit{\'e} au premier ordre : positionnement de cybergeo dans les cluster s{\'e}mantiques (un article pouvant {\^e}tre vu comme un vecteur de propbabilit{\'e}s d'appartenance aux disciplines)
\item Interdisciplinarit{\'e} au second ordre : liens de citation autour de cybergeo vers les diff{\'e}rentes disciplines
\item Comparaison des structures de communaut{\'e} : coefficient de clustering inter-couches ; donne tendance g{\'e}n{\'e}rale
\end{itemize}

}



\end{document}
